\documentclass[12pt]{article}
\usepackage[a4paper,margin=2.2cm]{geometry}
\usepackage{amsmath,amssymb,amsthm,mathtools}
\usepackage{hyperref}
\usepackage{cleveref}

% 中文（xelatex）
\usepackage{fontspec}
\usepackage{xeCJK}

\usepackage{fontspec}
\usepackage{xeCJK}
\setCJKmainfont{Noto Serif CJK SC}

\title{$A^{T}Ac = A^{T}u$}
\date{}
\begin{document}
\maketitle

\begin{flushleft}
    \subsubsection{非正交基下的真实投影}
    给定: \\
    $a_{1} = [1 ,1 ,0]^{T}, a_{2} = [1, 1, 1]^{T}, u = [4, 2, 2]^{T}$ \\
    令: \\
    $W = span(a_{1},a_{2})$, $a_{1},a_{2}$ 非正交

    思考问题: span 出来的的 W 是几维的? \\
    对于向量 $a_{1},a_{2}$ 它们的本质是没什么多大的区别的 是一个不同占比带方向的移动指令 \\
    span 出来的维度不与向量的行数量相关 因为向量的行数是分量的占比，向量终究只是一个有方向的箭头\\
    例如 $a_{2} = [1, 1, 1]^{T}$, 它的意思很简单: 这个指令箭头 是由 1份 x 基向量 + 1份 y 基向量 + 1份 z 基向量
    组成, 但它依然只是一个有方向的箭头, 所以 $W = span(a_{1},a_{2})$ 依然是2维平面 因为只有两个方向(这个两个方向线性无关), 但这个平面并不是
    标准的水平的 可能是倾斜的 它和$a_{1},a_{2}$ 的方向有直接关系

    1:计算 $a_{1} \cdot a_{2}$ \\
    \begin{align}
        a_{1} \cdot a_{2} \\
        = 1 * 1 + 1 * 1 + 0 * 1
        = 2
    \end{align}
    $a_{1},a_{2}$ 是否线性无关 ? \\
    是的, 它们线性无关 因为 $a_{1} \neq k \cdot a_{2} (k \in \mathbb{R})$ \\
    $a_{1},a_{2}$ 是否正交 ? \\
    不正交 $<a_{1},a_{2}> \neq  0 $ \\
    $a_{1},a_{2}$ 是否是 W 的一组非正交基 ? \\
    因为 $a_{1},a_{2}$线性无关 所以 $span(a_{1},a_{2})$ 是成立的, 是 W 的一组非正交基 \\

    2: 投影与残差向量 \\
    存在 $p = proj_{w}(u)$
    \begin{align}
        p = c_{1}a_{1} + c_{2}a_{2} \; ((c_{1},c_{2}) \in \mathbb{R}) \\
        = c_{1}[1 ,1 ,0]^{T} + c_{2}[1, 1, 1]^{T}                     \\
        =[c_{1} + c_{2}, c_{1} + c_{2}, c_{2}]^{T}                    \\
        r = u - p                                                     \\
        r = [4, 2, 2]^{T} - [c_{1} + c_{2}, c_{1} + c_{2}, c_{2}]^{T} \\
        r = [4 - c_{1}- c_{2}, 2 - c_{1}- c_{2}, 2 - c_{2}]^{T}       \\
    \end{align}
    已知 $r \perp W \rightarrow r \perp a_{1}, r \perp a_{2}$ \\
    那么我们可以使用内积算出 $c1,c2$ 存在下列代数式 \\
    \begin{align}
        c1 = (<r, a_{1}> = 0)                                                     \\
        = [4 - c_{1}- c_{2}, 2 - c_{1}- c_{2}, 2 - c_{2}]^{T} \cdot [1 ,1 ,0]^{T} \\
        = (4 - c_{1}- c_{2}) +  (2 - c_{1}- c_{2})  = 0                           \\
        =  6 - 2c_{1} -2c_{2} = 0                                                 \\
        =  2(3 - c_{1} - c_{2}) = 0                                               \\
        c2 = (<r, a_{2}> = 0)                                                     \\
        = [4 - c_{1}- c_{2}, 2 - c_{1}- c_{2}, 2 - c_{2}]^{T} \cdot [1 ,1 ,1]^{T} \\
        = (4 - c_{1}- c_{2}) +  (2 - c_{1}- c_{2})  + (2 - c_{2}) = 0             \\
        = 8 - 2c_{1} - 3c_{2} =  0                                                \\
        c_{2} = 2, c_{1} = 1
    \end{align}

    验证: \\
    $p = 1[1 ,1 ,0]^{T} + 2[1, 1, 1]^{T} = [3, 3, 2]^{T}$ \\
    $r = u - p = [4, 2, 2]^{T} - [3, 3, 2]^{T}  = [1, -1, 0]^{T} $ \\
    $<r, p> = 3 - 3 + 0 = 0 $ \\
    $<r,a_{1}> = 1 - 1 + 0 = 0$ \\
    $<r,a_{2}> = 1 - 1 + 0 = 0$ \\
    $u = p + r = [3, 3, 2]^{T} + [1, -1, 0]^{T} = [4, 2, 2]^{T}$ \\

    3:矩阵化 \\
    $span(a_{1}, a_{2})  \rightarrow A  = \begin{bmatrix}
            1 & 1 \\
            1 & 1 \\
            0 & 1 \\
        \end{bmatrix}$

    思考问题: 为什么 $span$ 可以矩阵化?  \\
    张成的本质是: 对于张成的子空间中的任何一个点 都可以通过列向量组合到达 (几何路线) \\
    矩阵: 这个矩阵的所有列向量能线性组合出的所有结果 (代数组合) \\
    对同一事物的两个视角 \\

    因为 矩阵是代数工具  \\
    那么第一个问题是: 是否存在一个标量向量 可以通过 A 矩阵映射到 u ? \\

    $c_{ideal} = \begin{bmatrix}
            c_{i1},
            c_{i2}
        \end{bmatrix}$ \\
    构建下列代数式: \\
    \begin{align}
        Ac_{ideal} = u \\
        \begin{bmatrix}
            1 & 1 \\
            1 & 1 \\
            0 & 1 \\
        \end{bmatrix} \begin{bmatrix}
                          c_{i1}, \\
                          c_{i2}
                      \end{bmatrix} =  \begin{bmatrix}
                                           4 \\
                                           2 \\
                                           2
                                       \end{bmatrix}
    \end{align}
    构建方程组: \\
    \begin{align}
        c_{i1} + c_{i2} = 4 \\
        c_{i1} + c_{i2} = 2 \\
        c_{i2} = 2          \\
    \end{align} \\
    我们发现方程组冲突, 无解 $Ac_{ideal} = u$. \\

    因为无解, 那么问题来了 我们能不能一个近似最佳解? 也就是 $A\hat{c} \approx u $\\

    这时代数和几何就发生了非常有意思的碰撞准确的来说是: transform \\
    在 W 子空间中 离 u 最近的点(几何) = 是 $A\hat{c} \approx u $ \\
    这个点就是 p(u 在 W 上的投影) \\
    $\hat{c} = \begin{bmatrix}
            c_{1},
            c_{2}
        \end{bmatrix}$ \\

    那么我们现在有这几个关键对象: $A, A^{T},\hat{c}, p = A\hat{c}, r = (u - p), u = (p + r) $ \\

    这里我们发现 r 有一个重要特性 $\to r \perp W  \land  r \perp a_{1} \land r \perp a_{2}$ \\
    也就是 $A^{T} r  = [0,0]^{T}$ 这里的几何意义很明显 r 没有 非零的 W 子空间方向的分量, 代数意义是 r 本来属于 $A^{T}$ 的零空间, 在映射转换过程中
    被抹去的 \\

    存在下列代数式
    \begin{align}
        A^{T}r = 0          \\
        A^{T}(u - p) = 0    \\
        A^{T}u - A^{T}p = 0 \\
        A^{T}u =  A^{T}p    \\
        A^{T}u = A^{T}A\hat{c}
    \end{align}

    我们需要仔细的分析上述代数式, 存在下列基础性质: \\
    1: 矩阵运算是一个线性变化 $A^{T}r$ 会得到一个新的向量
    2: $A^{T}r$ = $\begin{bmatrix}
            <a_{1},r> \\
            <a_{2},r>
        \end{bmatrix} = 0$ (矩阵乘法是行乘列) \\
    那么基于 1 的概念 $A^{T}r = r_{n}$ 因为 $r_{n} = 0$ 所以我们能确定 r 本身就是属于 $A^{T}$ 的零空间 \\
    3: $A^{T}u = \begin{bmatrix}
            <a_{1},u> \\
            <a_{2},u>
        \end{bmatrix} = u_{n}$ \\

    1: 这里我们来思考一个很有意思的点: $r_{n}, u_{n}$ 的空间关系? \\
    对于 $A^{T}/A$ 而言 它是一个矩阵, 也就是说它是一个线性变化的动作  \\
    $r$ 是一个属于$\mathbb{R}^{3}$, 那么$r_{n}$ 本质就是一个向量 \\
    我们观察一个向量所处于什么子空间, 关键是看两点: 向量的行元素量 和 向量的所处于的基向量组 \\
    $r_{n}$ 是一个 2d 向量, 那么它处于的基向量组是什么呢 ? \\
    我们知道 $r_{n}$ 的标量元素是来至 r 和 a 的内积, 那么我们能确定 $r_{n}$ 所处的基向量组是 ${a_{1}, a_{2}}$
    并不能 因为本质  $A^{T}/A$ 是一个变化动作 而非子空间, 这也是投影和矩阵变化的一个区别点
    投影是有目的地子空间的 也就是说是默认有一组基向量的
    但矩阵变化没有 它只是一种运算/映射动作??

    也就是说 $r_{n}$ 是一个2d向量(在2维空间, 默认是 e1,e2 的基向量组), 虽然 W 子空间也是一个2d平面的子空间 但 W 的基底是  ${a_{1}, a_{2}}$
    两个3维的基向量,两个基组的维度不一样 必然是不在一个空间的, 甚至 $r_{n}$ 不在 $\mathbb{R}^{3}$ 空间!


    2: $A^{T}$ 与 $A$  \\
    矩阵 它的本质是线性映射, 对于任意一个向量 它的默认基组是${e_{1}, ... , e_{k}}$ 这个基组是拥有维度的概念的也就是 Rank (线性无关的列数)\\
    对于矩阵而言, 它也有维度这个概念, 而且是两个 一个输入(配对矩阵乘法), 一个是输出(矩阵的列空间) \\
    A: \\
    $A \in \mathbb{R}^{3 \times 2}$ 它的输入维度 $\mathbb{R}^{2}$  输出维度是 $\mathbb{R}^{3}$ \\
    已知 $\hat{c} \in \mathbb{R}^{2 \times 1} \to A\hat{c} \in \mathbb{R}^{3 \times 1}$ \\
    这里有一个很重要的细节 虽然 $A\hat{c} \in \mathbb{R}^{3 \times 1} = \hat{c}_{new}$,
    $\hat{c}_{new}$ 到底处于是几维的? 这是一个非常重要的问题 \\
    对于代数形式结果来说  $A\hat{c} \in \mathbb{R}^{3 \times 1}$  它是有三个分量的 x, y, z  所以它处于的环境的确是 $\mathbb{R}^{3 \times 1}$ \\
    但问题是 A 的有效方向只有两个(2列) \\
    \begin{align}
        [a_{1}, a_{2}](A) \cdot [c_{1},c_{2}]^{T}(\hat{c}) \;\; ((a_{1}, a_{2} \in \mathbb{R}^{3}),(c_{1}, c_{2}) \in \mathbb{R})  \\
        =[a_{1}[1] * c_{1} + a_{2}[1] * c_{2}, \; a_{1}[2] * c_{1} + a_{2}[2] * c_{2}, \; a_{1}[3] * c_{1} + a_{2}[3] * c_{2}]^{T} \\
    \end{align}
    所以上面的代数证明了 $\hat{c}_{new} \in \mathbb{R}^{3 \times 1}$
    我们现在把 A 泛化 也就是说 $a_{1}[1], a_{1}[2],a_{2}[1],a_{2}[2],a_{1}[3],a_{2}[3]$ 完全是随机数并且是线性无关的, 但只有两列(两个方向) \\
    现在设: $[a_{1}[1] * c_{1} + a_{2}[1] * c_{2}, \; a_{1}[2] * c_{1} + a_{2}[2] * c_{2}, \; a_{1}[3] * c_{1} + a_{2}[3] * c_{2}]^{T} = [x, y, z]^{T}$ \\
    因为是线性无关的 那么必然存在关键代数关系 $xA + yB + zC = 0$ \\
    \begin{align}
        xA + yB + zC = 0                                                                                                                      \\
        (a_{1}[1] * c_{1} + a_{2}[1] * c_{2})A + (a_{1}[2] * c_{1} + a_{2}[2] * c_{2})B  + (a_{1}[2] * c_{1} + a_{2}[2] * c_{2})C  = 0        \\
        A(a_{1}[1] * c_{1}) + A(a_{2}[1] * c_{2}) + B(a_{1}[2] * c_{1}) + B(a_{2}[2] * c_{2}) + C(a_{1}[2] * c_{1}) + C(a_{2}[2] * c_{2}) = 0 \\
        c1(Aa_{1}[1] + Ba_{1}[2] + Ca_{1}[2]) +  c2(Aa_{2}[1]  +  Ba_{2}[2] + Ca_{2}[2]) = 0                                                  \\
        \text{得到关键方程组:}                                                                                                                       \\
        A a_{1}[1] + B a_{1}[2] + C a_{1}[3] = 0                                                                                              \\
        A a_{2}[1] + B a_{2}[2] + C a_{2}[3] = 0
    \end{align}
    这是一个关于未知数 A, B, C 的方程组（有 3 个未知数，却只有 2 个方程）。根据线性代数的定理，这样的方程组一定存在非零解！ 也就是说，我们一定能找到一组不全为零的常数 (A, B, C) 满足这两个方程。\\
    那么我们可以得到方程组: \\
    \begin{align}
        A + B + C = 0 ((A, B, C) \neq (0, 0, 0)) \\
        A + B = - C                              \\
        -C + C = 0
    \end{align}
    由 3列 退化成 2列 这就是 约束空间 它由 A 的有效方向确定 \\

    由上面的证明 我们可以解释: A 并不是一个简单的升维 或者说 A 的确会把结果映射到高维上 但这个结果真正处于的空间 是由 A 的 自由度(列)决定的

    $A^{T}$ 也是同理 它也会向另一个维度映射 但真正的约束空间是  $A^{T}$ 的 自由度(列)决定的

    对于矩阵运算来说 跨维度映射是非常重要的特性

    3: $A^{T}A$ \\
    $A^{T}A \in \mathbb{R}^{2 \times 2}$ \\
    \begin{align}
        A^{T}A \\
        = \begin{bmatrix}
              a_{1} \\
              a_{2}
          \end{bmatrix} \cdot \begin{bmatrix}
                                  a_{1} & a_{2}
                              \end{bmatrix}
        = \begin{bmatrix}
              a_{1} \cdot a_{1} & a_{1} \cdot a_{2}  \\
              a_{2} \cdot a_{1} & a_{2}  \cdot a_{2}
          \end{bmatrix}
    \end{align}

    从$A^{T}A = A_{n}$ 的输入角度来看 输入是 A 输出是 $A_{n} \in \mathbb{R}^{2 \times 2}$ \\
    这是一个非常明显的降维映射 目的地是2d \\
    如果 $<a_{1}, a_{2}> = 0 $ \\
    那么 $A_{n}$ 是对角矩阵 \\

    那么 $A^{T}A\hat{c}$ 几何意义是什么呢 ? \\
    输入一个 2维向量  $ \hat{c} \to A\hat{c}$ 升到 3维(投影) $ \to A^{T}$ 降到 2维(观测空间或对比空间?) \\

    上面的证明为了解决一个问题: 使用矩阵代数快速的找到近似最佳解(投影, 前题是 $Ac_{ideal} \neq u$ 无解) \\
    代数式表达为: $A^{T}A\hat{c} = A^{T}u$

\end{flushleft}


\end{document}